% GENERATED FILE. Edit canonical lessons, then run npm run generate:books.
% canonical-source-hash: fnv1a64:e6c1b841c8348010
% canonical-lessons: SA-C302-jalaughah, SA-C302-plavayati, SA-C302-pravahah, SA-C302-nimajjati, SA-C302-kramasah

\chapter{The Water Rose}
\label{ch:sa-the-water-rose}
\hlchaptermodality{302}

\begin{chapteropening}
\textbf{By the end of this chapter:} \emph{I can say how a flood rose little by little.}

Complete the last lesson of chapter 302: I can say how a flood rose little by little.
\end{chapteropening}

\section[\texorpdfstring{\sk{जलौघः} (jalaughaḥ)}{जलौघः (jalaughaḥ)}]{\sk{जलौघः} (jalaughaḥ) — a flood}

\label{lesson:SA-C302-jalaughah}



\begin{quote}
Before the new one: say the Sanskrit for flashes, then the Sanskrit for uproots.
\end{quote}

\subsection*{You'll want to know: \sk{जलौघः}}
\textbf{\sk{जलौघः}} — \emph{jalaughaḥ} — \textquotedblleft{}a flood\textquotedblright{}.

\begin{quote}
\sk{वर्षायाः} \sk{अनन्तरं} \sk{ग्रामे} \sk{जलौघः} \sk{आगच्छति।}
\end{quote}

\begin{quote}
\emph{varṣāyāḥ anantaraṃ grāme jalaughaḥ āgacchati.}
\end{quote}

\begin{quote}
After the rain a flood comes to the village.
\end{quote}

A masculine noun.

\begin{cousinweb}
\textbf{\sk{जलौघः}} (\emph{jalaughaḥ}) is built from \emph{jala}, \textquotedblleft{}water\textquotedblright{}, and \emph{ogha}, \textquotedblleft{}a flood, a rapid flow of water\textquotedblright{}; MW gives \emph{jalaugha} as a mass of water (MW).
\end{cousinweb}

\begin{grammarlens}[title={What you've built}]
One more word for telling what happened.
\end{grammarlens}

\subsection*{Your turn}
\begin{itemize}
\raggedright
  \item \emph{Say it:} \emph{jalaughaḥ}
  \item \emph{Say it:} \emph{jalaughaḥ}, once more
  \item \emph{Say it:} \emph{varṣāyāḥ anantaraṃ grāme jalaughaḥ āgacchati.} — After the rain a flood comes to the village.
  \item \emph{Recall:} say the Sanskrit for flashes, then the Sanskrit for uproots, then say \emph{jalaughaḥ} again
\end{itemize}

\begin{tcolorbox}[breakable,colback=teal!4,colframe=teal!35!black,title={Before you move on}]
What does \sk{जलौघः} mean? (\textquotedblleft{}A flood\textquotedblright{}.) Say it once more.
\end{tcolorbox}
\section[\texorpdfstring{\sk{प्लावयति} (plāvayati)}{प्लावयति (plāvayati)}]{\sk{प्लावयति} (plāvayati) — floods}

\label{lesson:SA-C302-plavayati}



\begin{quote}
Before the new one: say the Sanskrit for uproots, then the Sanskrit for a flood.
\end{quote}

\subsection*{You'll want to know: \sk{प्लावयति}}
\textbf{\sk{प्लावयति}} — \emph{plāvayati} — \textquotedblleft{}floods\textquotedblright{}.

\begin{quote}
\sk{नदी} \sk{क्षेत्राणि} \sk{प्लावयति।}
\end{quote}

\begin{quote}
\emph{nadī kṣetrāṇi plāvayati.}
\end{quote}

\begin{quote}
The river floods the fields.
\end{quote}

Past (it …): \sk{नदी} \sk{क्षेत्राणि} \sk{प्लावितवती}, \emph{nadī kṣetrāṇi plāvitavatī}.

\begin{cousinweb}
\textbf{\sk{प्लावयति}} (\emph{plāvayati}) comes from the root √plu, \textquotedblleft{}to float, swim\textquotedblright{}, which also gives \emph{plavate}, \textquotedblleft{}floats\textquotedblright{}; MW gives its causative \emph{plāvayati} as to make float, wash and inundate (MW).
\end{cousinweb}

\begin{grammarlens}[title={What you've built}]
One more word for telling what happened.
\end{grammarlens}

\subsection*{Your turn}
\begin{itemize}
\raggedright
  \item \emph{Say it:} \emph{plāvayati}
  \item \emph{Say it:} \emph{plāvayati}, once more
  \item \emph{Say it:} \emph{nadī kṣetrāṇi plāvayati.} — The river floods the fields.
  \item \emph{Recall:} say the Sanskrit for uproots, then the Sanskrit for a flood, then say \emph{plāvayati} again
\end{itemize}

\begin{tcolorbox}[breakable,colback=teal!4,colframe=teal!35!black,title={Before you move on}]
What does \sk{प्लावयति} mean? (\textquotedblleft{}Floods\textquotedblright{}.) Say it once more.
\end{tcolorbox}
\section[\texorpdfstring{\sk{प्रवाहः} (pravāhaḥ)}{प्रवाहः (pravāhaḥ)}]{\sk{प्रवाहः} (pravāhaḥ) — a current}

\label{lesson:SA-C302-pravahah}



\begin{quote}
Before the new one: say the Sanskrit for a flood, then the Sanskrit for floods.
\end{quote}

\subsection*{You'll want to know: \sk{प्रवाहः}}
\textbf{\sk{प्रवाहः}} — \emph{pravāhaḥ} — \textquotedblleft{}a current\textquotedblright{}.

\begin{quote}
\sk{नद्याः} \sk{प्रवाहः} \sk{वेगेन} \sk{वहति।}
\end{quote}

\begin{quote}
\emph{nadyāḥ pravāhaḥ vegena vahati.}
\end{quote}

\begin{quote}
The river's current flows fast.
\end{quote}

A masculine noun.

\begin{cousinweb}
\textbf{\sk{प्रवाहः}} (\emph{pravāhaḥ}) is built from \emph{pra}, \textquotedblleft{}forth\textquotedblright{}, and the root √vah, \textquotedblleft{}to carry, flow\textquotedblright{}, as in \emph{pravahati}, \textquotedblleft{}flows\textquotedblright{}; MW gives \emph{pravāha} as a stream or current (MW).
\end{cousinweb}

\begin{grammarlens}[title={What you've built}]
One more word for telling what happened.
\end{grammarlens}

\subsection*{Your turn}
\begin{itemize}
\raggedright
  \item \emph{Say it:} \emph{pravāhaḥ}
  \item \emph{Say it:} \emph{pravāhaḥ}, once more
  \item \emph{Say it:} \emph{nadyāḥ pravāhaḥ vegena vahati.} — The river's current flows fast.
  \item \emph{Recall:} say the Sanskrit for a flood, then the Sanskrit for floods, then say \emph{pravāhaḥ} again
\end{itemize}

\begin{tcolorbox}[breakable,colback=teal!4,colframe=teal!35!black,title={Before you move on}]
What does \sk{प्रवाहः} mean? (\textquotedblleft{}A current\textquotedblright{}.) Say it once more.
\end{tcolorbox}
\section[\texorpdfstring{\sk{निमज्जति} (nimajjati)}{निमज्जति (nimajjati)}]{\sk{निमज्जति} (nimajjati) — sinks; goes under}

\label{lesson:SA-C302-nimajjati}



\begin{quote}
Before the new one: say the Sanskrit for floods, then the Sanskrit for a current.
\end{quote}

\subsection*{You'll want to know: \sk{निमज्जति}}
\textbf{\sk{निमज्जति}} — \emph{nimajjati} — \textquotedblleft{}sinks; goes under\textquotedblright{}.

\begin{quote}
\sk{जले} \sk{नौका} \sk{शनैः} \sk{निमज्जति।}
\end{quote}

\begin{quote}
\emph{jale naukā śanaiḥ nimajjati.}
\end{quote}

\begin{quote}
The boat slowly goes under in the water.
\end{quote}

Past (it …): \sk{नौका} \sk{जले} \sk{निमग्नवती}, \emph{naukā jale nimagnavatī}.

\begin{cousinweb}
\textbf{\sk{निमज्जति}} (\emph{nimajjati}) is built from \emph{ni}, \textquotedblleft{}down\textquotedblright{}, and the root √majj, \textquotedblleft{}to sink, dive\textquotedblright{}, as in \emph{majjati}; MW gives \emph{ni-majj} as to sink down, plunge into or disappear (MW).
\end{cousinweb}

\begin{grammarlens}[title={What you've built}]
One more word for telling what happened.
\end{grammarlens}

\subsection*{Your turn}
\begin{itemize}
\raggedright
  \item \emph{Say it:} \emph{nimajjati}
  \item \emph{Say it:} \emph{nimajjati}, once more
  \item \emph{Say it:} \emph{jale naukā śanaiḥ nimajjati.} — The boat slowly goes under in the water.
  \item \emph{Recall:} say the Sanskrit for floods, then the Sanskrit for a current, then say \emph{nimajjati} again
\end{itemize}

\begin{tcolorbox}[breakable,colback=teal!4,colframe=teal!35!black,title={Before you move on}]
What does \sk{निमज्जति} mean? (\textquotedblleft{}Sinks; goes under\textquotedblright{}.) Say it once more.
\end{tcolorbox}
\section[\texorpdfstring{\sk{क्रमशः} (kramaśaḥ)}{क्रमशः (kramaśaḥ)}]{\sk{क्रमशः} (kramaśaḥ) — by degrees; step by step}

\label{lesson:SA-C302-kramasah}



\begin{quote}
Before the new one: say the Sanskrit for a current, then the Sanskrit for goes under.
\end{quote}

\subsection*{You'll want to know: \sk{क्रमशः}}
\textbf{\sk{क्रमशः}} — \emph{kramaśaḥ} — \textquotedblleft{}by degrees; step by step\textquotedblright{}.

\begin{quote}
\sk{जलं} \sk{क्रमशः} \sk{वर्धते।}
\end{quote}

\begin{quote}
\emph{jalaṃ kramaśaḥ vardhate.}
\end{quote}

\begin{quote}
The water rises by degrees.
\end{quote}

\begin{cousinweb}
\textbf{\sk{क्रमशः}} (\emph{kramaśaḥ}) is built on \emph{krama}, \textquotedblleft{}a step, an order\textquotedblright{}, with the ending \emph{-śas}; MW gives \emph{kramaśas} as gradually, by degrees (MW).
\end{cousinweb}

\begin{grammarlens}[title={What you've built}]
One more word for telling what happened.
\end{grammarlens}

\subsection*{Your turn}
\begin{itemize}
\raggedright
  \item \emph{Say it:} \emph{kramaśaḥ}
  \item \emph{Say it:} \emph{kramaśaḥ}, once more
  \item \emph{Say it:} \emph{jalaṃ kramaśaḥ vardhate.} — The water rises by degrees.
  \item \emph{Recall:} say the Sanskrit for a current, then the Sanskrit for goes under, then say \emph{kramaśaḥ} again
\end{itemize}

\begin{tcolorbox}[breakable,colback=teal!4,colframe=teal!35!black,title={Before you move on}]
What does \sk{क्रमशः} mean? (\textquotedblleft{}By degrees; step by step\textquotedblright{}.) Say it once more.
\end{tcolorbox}
